\documentclass[10pt,a4paper]{amsart}
\usepackage[margin=19mm]{geometry}
\usepackage{amsmath,amssymb,mathtools}
\usepackage[scr=rsfs,cal=euler]{mathalfa}
\usepackage{xfrac}
\usepackage[unicode,hidelinks,bookmarks=false]{hyperref}
\newcommand{\ie}{i.e.\ }
\newcommand{\cf}{cf.\ }
\newcommand{\resp}{respectively\ }
\newcommand{\Subsection}{\subsection}
\providecommand{\nobreakdash}{\nobreak\hbox{-}}
\setlength{\emergencystretch}{2em}
\multlinegap0pt
\usepackage{bm}
\usepackage{tikz}
\usepackage[normalem]{ulem}
\usepackage{multirow}
\usepackage{natbib}
\usepackage{float}
\usepackage{amssymb}
\usepackage{xcolor}
\usepackage{siunitx}
\usepackage{booktabs}
\newtheorem{lemma}{Lemma}
\newcommand{\Reynolds}{{\operatorname{Re}}}
\newcommand{\bU}{{\mathbf{U}}}
\newcommand{\bomega}{{\boldsymbol{\omega}}}
\newcommand{\bu}{{\mathbf{u}}}
\newcommand{\bv}{{\mathbf{v}}}
\title{High-order DLM-ALE discretizations with robust operator preconditioning for fluid-rigid-body interaction}
\author{Qi Xin, Shihua Gong, Lingyue Shen, Pinjing Wen, Yumiao Zhang, Yan Chen, Jiarui Han, Jinchao Xu}
\date{Source: arXiv:2602.01094v2 (CC BY 4.0)}
\begin{document}
\maketitle

\noindent\emph{Source and licence.} High-order DLM-ALE discretizations with robust operator preconditioning for fluid-rigid-body interaction. Version arXiv:2602.01094v2 (CC BY 4.0), distributed by arXiv under the Creative Commons Attribution 4.0 International licence, http://creativecommons.org/licenses/by/4.0/, as recorded in the arXiv licence metadata for this exact version and shown on the abstract page https://arxiv.org/abs/2602.01094v2. Full licence notice: \texttt{licenses/arxiv-2602.01094-CC-BY-4.0.txt}.\par\smallskip

\noindent\textit{Editorial scope.} Counted unit: the article's lemma lemma (source0.tex lines 1281--1286). The article's complete original proof of it is delivered with it (source0.tex lines 1287--1381, one \texttt{proof} environment of 95 lines). No mathematics is written, repaired or re-derived: the statement and the proof are the article's own lines, in the article's own notation, and no retained line is substituted or re-flowed.  Retained uncounted as the source-local context the proof needs: lines 1101--1107.  Nothing was omitted from any retained span, so the package carries no omission record.  No retained line contains a citation, so the wrapper carries no bibliography and no bibliography entry.  No retained line contains a figure environment, an image-inclusion command or drawing code, so no image is delivered and the package declares no asset dependency.  Editorial formatting only, in addition: an AMS \texttt{amsart} preamble that restates the article's own declarations and macros used by the retained lines, namely the environments \texttt{lemma} on separate counters, together with the standard packages those lines need.  The retained environments print in this document as Lemma 1 in order of appearance, whereas the article prints the same items with the numbering of its own full document.  The complete native source of the article as distributed in the arXiv e-print for this exact version is bundled unchanged as \texttt{sources/193759/source0.tex}.
\begin{equation}\label{eq:tildeS}
    \begin{aligned}
        &M: Q \to Q', \quad \langle Mp, q \rangle = (p, q)_{\Omega_f^t}, \\
        &A: Q \to Q', \quad \langle Ap, q \rangle = (\nabla p, \nabla q)_{\Omega_f^t}, \\
        &\tilde{S} = (M^{-1} + \frac{\Reynolds}{\gamma\Delta t} A^{-1})^{-1}.
    \end{aligned}
\end{equation}

\begin{lemma}
    The continuity condition holds for the bilinear forms $b^t(\cdot,\cdot)$ and $c^t(\cdot,\cdot)$, i.e., there exists a positive constant $C$ independent of $\Reynolds$, $\gamma$ and $\Delta t$ such that for all $\bu\in V$, ${q}\in Q$, ${\boldsymbol{\mu}}\in\boldsymbol{\Lambda}$, the following estimate holds
    \begin{equation}
        {b}^t({q},\bu)+{c}^t({\boldsymbol{\mu}},\bu - {\bU} - {\bomega}\times\mathbf{r})\leq C(\|{q}\|_Q^2+\|{\boldsymbol{\mu}}\|_{\boldsymbol{\Lambda}}^2)^\frac12 (\|\bu\|_V^2 + \|{\bU}\|_{T_M}^2 + \|{\bomega}\|_{R_M}^2)^\frac12.
    \end{equation}
\end{lemma}

\begin{proof}
    The continuity condition for the bilinear form $c^t(\boldsymbol{\mu},\bu - \bU - \bomega\times\mathbf{r})$ is  straightforward. Indeed,
\begin{equation}
    \begin{aligned}
        {c}^t({\boldsymbol{\mu}},\bu - {\bU} - {\bomega}\times\mathbf{r})
        &\leq \|{\boldsymbol{\mu}}\|_{\boldsymbol{\Lambda}} \left( \|\bu\|_V + \|\bU\|_{T_M} + \|\bomega\|_{R_M} \right).
    \end{aligned}
\end{equation}
However, the continuity condition for the bilinear form $b^t(q,\bu)$ requires more detailed analysis. 
We define the linear operators $B: V \to Q'$, $M_\bv: V \to V'$, and $A_\bv: V \to V'$ by
\begin{equation}
    \begin{aligned}
        &\langle B\bu, q \rangle = b^t(q, \bu), \\
        &\langle A_\bv\bu, \bv \rangle = a^t(\bu, \bv), \\
        &\langle M_\bv\bu, \bv \rangle = m^t(\bu, \bv),
    \end{aligned}
\end{equation}
and let $B^*: Q \to V'$ be the adjoint of $B$, and $T_\bv = \frac{\Reynolds}{\gamma\Delta t} M_\bv + A_\bv$. Since $A_\bv$ and $M_\bv$ are self-adjoint and positive definite, we obtain the identity
\begin{equation}
    \begin{aligned}
    \left( \sup_{\bu \in V} \frac{b^t(q, \bu)}{\|\bu\|_V} \right)^2 
    &= \sup_{\bu \in V} \frac{\langle B^* q, \bu \rangle^2}{\langle T_\bv \bu, \bu \rangle} 
     = \sup_{\bu \in V} \frac{\langle T_\bv^{-1} B^* q, T_\bv \bu \rangle^2}{\langle T_\bv^{-1} T_\bv \bu, T_\bv \bu \rangle} 
     = \langle T_\bv^{-1} B^* q, B^* q \rangle 
     = \langle B T_\bv^{-1} B^* q, q \rangle.
    \end{aligned}
\end{equation}
Denoting $S = B T_\bv^{-1} B^*$, it suffices to show that
\begin{equation}
    \langle S q, q \rangle \leq C \langle \tilde{S} q, q \rangle = C \|q\|_Q^2
\end{equation}
for all $q \in Q$. This is equivalent to
\begin{equation}
    \langle S \tilde{S}^{-1} q', \tilde{S}^{-1} q' \rangle \leq C \langle \tilde{S}^{-1} q', q' \rangle
\end{equation}
for all $q' \in Q'$. It therefore suffices to show that the spectral radius of $S \tilde{S}^{-1}$ is bounded by $C$.
It is easy to show that $S \tilde{S}^{-1}$ and $\tilde{S}^{-1} S$ have the same spectrum.
As the closure of the numerical range of a self-adjoint operator with respect to an inner product defined by $\langle \cdot, S\cdot \rangle$ is the convex hull of its spectrum, it suffices to show that
\begin{equation}
    \langle \tilde{S}^{-1} S q, S q \rangle \leq C \langle S q, q \rangle
\end{equation}
for all $q \in Q$. By the definition of $\tilde{S}$ in \eqref{eq:tildeS}, we have
\begin{equation}
    \begin{aligned}
        \langle \tilde{S}^{-1} S q, S q \rangle = \frac{\Reynolds}{\gamma\Delta t} \langle A^{-1} S q, S q \rangle + \langle M^{-1} S q, S q \rangle.
    \end{aligned}
\end{equation}
For the first term, we have
\begin{equation}
    \begin{aligned}
        \frac{\Reynolds}{\gamma\Delta t} \langle A^{-1} S q, S q \rangle 
        &= \frac{\Reynolds}{\gamma\Delta t} \sup_{p \in Q} \frac{\langle S q, p \rangle^2}{\langle A p, p \rangle} 
         = \sup_{p \in Q} \frac{\langle T_\bv^{-1} B^* q, B^* p \rangle^2}{(\nabla p, \nabla p)_{\Omega_f^t}} 
         = \frac{\Reynolds}{\gamma\Delta t} \sup_{p \in Q} \frac{(T_\bv^{-1} B^* q, \nabla p)^2}{(\nabla p, \nabla p)_{\Omega_f^t}} \\
        &= \frac{\Reynolds}{\gamma\Delta t} \|T_\bv^{-1} B^* q\|_{L^2(\Omega_f^t)}^2 
         = \left\langle \frac{\Reynolds}{\gamma\Delta t} M_\bv T_\bv^{-1} B^* q, T_\bv^{-1} B^* q \right\rangle \\
        &\leq \left\langle \left( \frac{\Reynolds}{\gamma\Delta t} M_\bv + A_\bv \right) T_\bv^{-1} B^* q, T_\bv^{-1} B^* q \right\rangle 
         = \langle S q, q \rangle.
    \end{aligned}
\end{equation}
For the second term, we claim that
\begin{equation}
    \begin{aligned}
        \langle M^{-1} S q, S q \rangle \leq \langle S q, q \rangle.
    \end{aligned}
\end{equation}
To prove this, we again apply the argument on the spectral radius of a self-adjoint operator, so it suffices to show that
\begin{equation}
    \langle M^{-1} S q, M q \rangle \leq \langle M q, q \rangle,
\end{equation}
i.e.,
\begin{equation} 
    \langle S q, q \rangle \leq \langle M q, q \rangle.
\end{equation}
We have
\begin{equation}
    \begin{aligned}
        \langle S q, q \rangle 
        &= \langle T_\bv^{-1} B^* q, B^* q \rangle 
         = \langle T_\bv^{-1} B^* q, T_\bv T_\bv^{-1} B^* q \rangle \\
        &= \sup_{\bu \in V} \frac{\langle T_\bv^{-1} B^* q, T_\bv \bu \rangle^2}{\langle T_\bv \bu, \bu \rangle} 
         \leq C \sup_{\bu \in V} \frac{\langle B \bu, q \rangle^2}{\|\bu\|_{H^1(\Omega_f^t)}^2} \\
        &\leq C \langle M q, q \rangle,
    \end{aligned} 
\end{equation}
where we used $\|\nabla \cdot \bu| \leq C \|\bu\|_{H^1(\Omega_f^t)}$ and Korn's inequality. Combining both terms, we obtain
\begin{equation}
    \langle \tilde{S}^{-1} S q, S q \rangle \leq C \langle S q, q \rangle.
\end{equation}
This implies that the spectral radius of $S \tilde{S}^{-1}$ is bounded by $C$, and hence
\begin{equation}
    \sup_{\bu \in V} \frac{b^t(q, \bu)}{\|\bu\|_V} \leq C \|q\|_Q.
\end{equation}
Thus, the continuity condition for the bilinear form $b^t(q, \bu)$ is established.
\end{proof}

\end{document}
